% 2008 年广东卷文科第 11 题：频率分布直方图。
% 来源 PDF：gaokaomath-src/普通高考/2008/2008广东文.pdf
% SHA-256：1adedadb12b82b7b4d04721908bd686e747021593ec91da14b0e4e979929373c
% 原页：2008广东文.pdf 第 1 页（A3 三栏）第 2 栏，书页码 447。
% 重绘依据：横轴在 0 与 45 之间断开，组距 10；五组 [45,55)、[55,65)、[65,75)、[75,85)、
%   [85,95) 的“频率/组距”依次为 0.020、0.040、0.025、0.010、0.005，纵轴每 0.005 一格并配虚线。
\begin{tikzpicture}[x=1.0cm, y=148cm, semithick, font=\small,
  every node/.style={inner sep=1.2pt}]
  % 坐标轴
  \draw[->,>=stealth] (0,0) -- (12.0,0) node[right] {产品数量};
  \draw[->,>=stealth] (0,0) -- (0,0.0448) node[above] {$\dfrac{\text{频率}}{\text{组距}}$};
  \node[below left] at (0,0) {$0$};

  % 横轴断裂符号
  \draw (0.18,0) -- (0.30,0.0024) -- (0.42,0) -- (0.54,0.0024) -- (0.66,0);

  % 纵轴刻度与水平虚线
  \foreach \h/\lab in {%
    0.005/0.005, 0.010/0.010, 0.015/0.015, 0.020/0.020,
    0.025/0.025, 0.030/0.030, 0.035/0.035, 0.040/0.040} {
    \draw (0,\h) -- (-0.12,\h);
    \node[left] at (-0.12,\h) {\lab};
    \draw[dashed] (0,\h) -- (11.4,\h);
  }

  % 横轴刻度
  \foreach \x/\lab in {1/45, 3/55, 5/65, 7/75, 9/85, 11/95} {
    \draw (\x,0) -- (\x,-0.0024);
    \node[below] at (\x,-0.0024) {\lab};
  }

  % 各矩形（宽 10，高为频率/组距）
  \draw (1,0) rectangle (3,0.020);
  \draw (3,0) rectangle (5,0.040);
  \draw (5,0) rectangle (7,0.025);
  \draw (7,0) rectangle (9,0.010);
  \draw (9,0) rectangle (11,0.005);
\end{tikzpicture}
